% Ad Familiares 1.5a (ad-familiares-01-05a)
% Source: https://raw.githubusercontent.com/PerseusDL/canonical-latinLit/master/data/phi0474/phi056/phi0474.phi056.perseus-lat1.xml
% Fetched by scripts/fetch_latin.py

% Dateline (Perseus): Scr. Romae circiter Non. Febr. a. 698 (56) .

\ciceroLetterOpener{M. CICERO S.D. P. LENTVLO PROCOS.}

\ciceroSection{1} tametsi mihi nihil fuit optatius quam ut primum abs te ipso, deinde a ceteris omnibus quam gratissimus erga te esse cognoscerer tamen adficior summo dolore eius modi tempora post tuam profectionem consecuta esse, ut et meam et ceterorum erga te fidem et benevolentiam absens experirere; te videre et sentire eandem fidem esse hominum in tua dignitate, quam ego in mea salute sum expertus, ex tuis litteris intellexi.

\ciceroSection{2} nos cum maxime consilio, studio, labore, gratia de causa regia niteremur, subito exorta est nefaria Catonis promulgatio, quae nostra studia impediret et animos a minore cura ad summum timorem traduceret; sed tamen, in eius modi perturbatione rerum quamquam omnia sunt metuenda, nihil magis quam perfidiam timemus et Catoni quidem, quoquo modo se res habet, profecto resistemus.

\ciceroSection{3} de Alexandrina re causaque regia tantum habeo polliceri, me tibi absenti tuisque praesentibus cumulate satis facturum, sed vereor, ne aut eripiatur causa regia nobis aut deseratur; quorum utrum minus velim, non facile possum existimare. sed , si res coget, est quiddam tertium, quod neque Selicio nec mihi displicebat, ut neque iacere pateremur nec nobis repugnantibus ad eum deferri, ad quem prope iam delatum existimatur. A nobis agentur omnia diligenter, ut neque, si quid obtineri poterit, non contendamus nec, si quid non obtinuerimus, repulsi esse videamur;

\ciceroSection{4} tuae sapientiae magnitudinisque animi est omnem amplitudinem et dignitatem tuam in virtute atque in rebus gestis tuis atque in tua gravitate positam existimare; si quid ex iis rebus, quas tibi fortuna largita est, non nullorum hominum perfidia detraxerit, id maiori illis fraudi quam tibi futurum. A me nullum tempus praetermittitur de tuis rebus et agendi et cogitandi; utor ad omnia Q. Selicio, neque enim prudentiorem quemquam ex tuis neque fide maiore esse iudico neque amantiorem tui
